\section{Introduction}
\label{sec:introduction}

Multi-hop intelligent information retrieval requires complete evidence acquisition. Retrieval-augmented generation (RAG)~\cite{lewis2020rag} grounds answers in external sources, but a multi-hop question can require documents linked by entities, comparisons or temporal relations. Missing one source may break the evidence chain despite a high rank for another. We therefore use complete supporting-title recovery as the primary retrieval diagnostic, supplemented by exact processed-record and evidence-interface audits.

Adaptive systems change queries or retrieve repeatedly to recover missing evidence. Multi-hop dense retrieval~\cite{xiong2021multihopdense} learns sequential evidence acquisition, and IRCoT~\cite{trivedi2022ircot} interleaves retrieval with reasoning. ReAct~\cite{yao2022react}, Self-RAG~\cite{asai2023selfrag} and Adaptive-RAG~\cite{jeong2024adaptiverag} add acting, reflection or complexity-based control, respectively. A final score alone cannot distinguish query representation, generator knowledge, corpus effects and additional calls. RCT-MARS~\cite{srivastava2026rctmars} shows that learned retrieval routing can fail to beat a strong fixed strategy, motivating matched diagnostics before crediting adaptive control.

HyDE~\cite{gao2022hyde} generates a hypothetical passage containing likely entities and relations for retrieval. Here the question and passage form the query; the reader receives only the original question and retrieved corpus evidence. This boundary keeps query generation separate from answer evidence. Case-based and concept-based retrieval~\cite{nasar2024retrieval} illustrates the role of representation and similarity choices, while BPLLM~\cite{bernardi2024bpllm} motivates component-wise validation in modular RAG systems. We study expansion as one such component.

The evaluation problem is conditional transport: which conclusions remain supported after a system component changes? Hosted aliases may drift, generated passages may reproduce benchmark content, and candidate pools may conceal corpus-scale difficulty. Data-centric AI~\cite{malerba2024datacentric} and FactualRAG~\cite{chang2026factualrag} motivate separating data and RAG-system evaluation concerns. The unresolved question here is whether matched expansion comparisons remain favorable when retrieved support passes through the actual reader interface. We examine this chain using complete-support retrieval, retained annotated evidence and paired answer scores. Generator, retriever, corpus, serializer, reader and budget identify each tested configuration. Existing paired inference supplies uncertainty; the contribution comes from the observed boundaries and their consequences for component evaluation.

We study four research questions. \textbf{RQ1:} Does HyDE-style expansion improve multi-hop evidence acquisition in the tested controlled pipelines? \textbf{RQ2:} How do matched gains differ across the tested generators? \textbf{RQ3:} Do retrieval conclusions survive full-wiki scale and an independently developed answer reader? \textbf{RQ4:} What evidence/answer trade-off appears as IRCoT receives more retrieval rounds? The corresponding contributions are:
\begin{enumerate}
\item \textbf{C1. Expansion against matched question-only baselines.} Controlled generator comparisons distinguish the hosted artifact from pinned open-generator estimates. News comparisons against the same question-only Dense retriever identify negative standard-expansion retrieval and answer effects. This tests expansion beyond comparison with a lexical baseline.
\item \textbf{C2. Retrieved versus delivered evidence.} Character and final-token audits locate losses that retrieval scores cannot reveal. A matched budget intervention increases annotated-support retention without establishing answer gains, separating delivery diagnostics from answer evaluation.
\item \textbf{C3. Retrieval versus answer gains.} Corpus-scale cross-encoder comparisons and two pinned readers identify configurations with supported retrieval gains but unresolved answer gains. Reader controls provide a separate capability reference, while round prefixes describe iteration under unequal accumulated cost.
\end{enumerate}

Here, reliability boundaries denote the empirical limits of supported retrieval and answer gains within the tested configurations. An unresolved gain is not evidence of no effect. Named contrasts, family-specific paired decisions and hashed records support inspection; they do not yield a universal account of when expansion works, a predictive model or a validated method-selection rule.

The tested English configurations comprise a complete HotpotQA Wikipedia snapshot, a 2Wiki official split-union corpus and MultiHop-RAG~\cite{tang2024multihoprag} news documents. The latter's local IRCoT condition is unavailable, and no independent human evidence-sufficiency validation was performed. Available conditions do not complete the registered transport test. The remainder of the article presents related work (Section~\ref{sec:related_work}), the protocol (Section~\ref{sec:method}), RQ1--RQ4 and external validation (Section~\ref{sec:experiments}), and interpretation and limits (Sections~\ref{sec:discussion_limitations}--\ref{sec:conclusion}).
